% Part: incompleteness
% Chapter: representability-in-q
% Section: c

\documentclass[../../../include/open-logic-section]{subfiles}

\begin{document}

\olfileid{inc}{req}{cee}
\olsection{Fungsi-fungsi $C$}

Anggep $C$ minangka himpunan fungsi paling cilik kang ngemot
\begin{enumerate}
\item $0$,
\item fungsi penerus,
\item panjumlahan,
\item perkalian,
\item proyeksi, lan
\item fungsi karakteristik kanggo kesamaan, $\Char{=}$;
\end{enumerate}
lan katutup tumrap
\begin{enumerate}
\item komposisi, lan
\item telusuran tanpa wates kang ditrapaké marang
fungsi reguler.
\end{enumerate}
Élinga yèn watesan pungkasan iki mung ateges operasi
$\mu$ bisa digunakaké yèn asilé total. Bandhingna karo
definisi fungsi \emph{rekursif umum}: ing kéné kita wis
nambahaké panjumlahan, perkalian, lan $\Char{=}$,
nanging ora nyakup rekursi primitif.

Cetha yèn saben fungsi ing $C$ rekursif, amarga panjumlahan,
perkalian, lan $\Char{=}$ uga rekursif. Kita bakal nuduhaké
yèn walikané uga bener; tegesé, kanthi unsur-unsur liya
ing $C$, kita bisa nindakake rekursi primitif.

\end{document}
